\section{Independent calibration and prediction
error}\label{calibration-independent-calibration-and-prediction-error}

Part of \href{DRAFT.md}{the manuscript}. E-numbering and section numbers
are continuous across \href{scalar_exit.md}{scalar exit},
\href{calibration.md}{calibration}, and \href{risk_boundaries.md}{risk
boundaries}. Full statements and proofs are retained. Constants are not
numerically calibrated; integration is not a completed independent proof
audit. Use \(B_U=w/4\) for the update bound and \(M\) for calibration
sample size; the original Gaussian sections use \(M=w/4\) locally.

\subsection{3. Quantitative stability of the
geometry}\label{calibration-quantitative-stability-of-the-geometry}

Let \(\widetilde v\) be another entropy field, or a smooth perturbation
whose map \(\widetilde G=x+\eta\widetilde v\) is an increasing
bijection. Compare it with the fixed field \(v\) on \(\mathcal K\). All
derivatives through order three are assumed uniformly bounded in the
neighborhood of fields under consideration.

\subsubsection{Lemma E.4. Stability of the root and normalized
conjugacy}\label{calibration-lemma-e.4.-stability-of-the-root-and-normalized-conjugacy}

There are constants \(d_*,C>0\) such that, if \[
d_v:=\|\widetilde v-v\|_{C^3(\mathcal K)}\le d_*,
\tag{E.33}
\] then \(\widetilde v\) has a unique zero \(\widetilde q\) on
\(\mathcal K\), has the same outward sign pattern relative to
\(\widetilde q\), and \[
|\widetilde q-q|\le C d_v,\qquad
|\widetilde\lambda-\lambda|+|\widetilde A-A|\le C d_v.
\tag{E.34}
\] Here \(\widetilde\lambda=\widetilde v'(\widetilde q)\) and
\(\widetilde A=1+\eta\widetilde\lambda\).

The normalized conjugacies satisfy \[
\|\widetilde H-H\|_{C^1(\mathcal K)}\le C d_v,
\qquad
\|\widetilde H\|_{C^2(\mathcal K)}\le C.
\tag{E.35}
\] For the critical initialization in G4, writing
\(d_q=\widetilde q-q\), \[
\sqrt n\{\widetilde H(x_{0,n})-H(x_{0,n})\}
=-\sqrt n\,d_q+
O\!\left(d_v+\sqrt n\,d_q^2\right).
\tag{E.36}
\]

\textbf{Proof.} Choose a local interval about \(q\) on which
\(v'\ge\lambda/2\). A sufficiently small \(C^1\) perturbation has
positive derivative there and opposite signs at its endpoints. It
therefore has exactly one local zero. Away from this interval,
compactness and G1 give a strictly positive lower bound on \(|v|\), so
no additional zeros appear under a sufficiently small perturbation. The
mean-value theorem gives the root bound; differentiation and bounded
\(v''\) give the curvature and multiplier bounds.

For completeness, the quantitative dependence of \(H\) can be proved
directly from the inverse construction. Center the inverse maps at their
own roots, writing them as \(g,\widetilde g\). On a common sufficiently
small local interval, \[
\|g-\widetilde g\|_{C^2}\le C d_v,\qquad
|g'|,|\widetilde g'|\le b<1,\qquad
\overline A b^2<1,
\tag{E.37}
\] where \(\overline A\) bounds both multipliers. The inverse-function
formulas and the common lower bound on the forward derivatives justify
the first inequality. Centering introduces only \(O(d_v)\) changes,
including in the second derivatives, by the \(C^3\) hypothesis.

Both inverse maps fix zero. Iterating the difference recursion and its
first derivative yields \[
|\widetilde g^j(y)-g^j(y)|
\le C d_v(j+1)b^{j-1}|y|,
\] \[
|(\widetilde g^j)'(y)-(g^j)'(y)|
\le C d_v(j+1)^2b^{j-1}.
\tag{E.38}
\] For the derivative estimate, use the product rule, the common
second-derivative bound, and the first estimate. The harmless cases
\(j=0,1\) are absorbed in \(C\).

Let \(R_g(y)=g(y)-A^{-1}y\), with the analogous definition for
\(\widetilde g\). Their difference has value and derivative zero at zero
and \(C^2\) norm at most \(C d_v\). Thus \[
|R_{\widetilde g}(y)-R_g(y)|\le C d_v y^2,\qquad
|R_{\widetilde g}'(y)-R_g'(y)|\le C d_v|y|.
\] Use the convergent series \[
h(y)=y+\sum_{j\ge0}A^{j+1}R_g(g^j(y)),
\tag{E.39}
\] and its perturbed version. Equations (E.38), the quadratic vanishing
of the residuals, and \[
|\widetilde A^{j+1}-A^{j+1}|
\le C d_v(j+1)\overline A^j
\] bound the \(C^1\) difference of the \(j\)-th terms by \[
C d_v(j+1)^3(\overline A b^2)^j.
\] Summation proves local \(C^1\) Lipschitz dependence. The unperturbed
construction of Lemma E.2, with uniform bounds, also proves a common
\(C^2\) bound for all the local conjugacies.

A common finite number of inverse steps maps \(\mathcal K\) into the
local domain for every sufficiently close field. To see uniformity,
outside the local domain the outward drift is bounded away from zero,
while inside it the inverse is uniformly contracting. The extension
formula (E.13), finite composition, and the inverse derivative bounds
give (E.35) on \(\mathcal K\).

Finally, since \(\widetilde H(\widetilde q)=0\) and
\(\widetilde H'(\widetilde q)=1\), \[
\widetilde H(q)=-d_q+O(d_q^2).
\] The integral of \(\widetilde H'-H'\) between \(q\) and \(x_{0,n}\) is
at most \(C d_v/\sqrt n\). Multiplication by \(\sqrt n\) proves (E.36).
\(\square\)

This lemma is local in the space of fields. Its constants can be large
if the inverse orbit spends a long time away from the root, if a
boundary approaches another equilibrium, or if \(\lambda\) is small. It
does not assert stable prediction across a change in the number or type
of equilibria.

\subsection{4. Independent unlabeled field
estimation}\label{calibration-independent-unlabeled-field-estimation}

Use the sample (E.3) to define \[
\widehat v_M(x)=\frac1M\sum_{i=1}^M U(x,\widetilde Z_i).
\tag{E.40}
\] On the event that it has the corresponding root \(\widehat q_M\) on
\(\mathcal K\), set \[
\widehat\lambda_M=\widehat v_M'(\widehat q_M),\qquad
\widehat A_M=1+\eta\widehat\lambda_M,
\] \[
\widehat\sigma_{g,M}^2
=\frac1M\sum_{i=1}^M U(\widehat q_M,\widetilde Z_i)^2,
\qquad
\widehat s_M^2
=\frac{\eta^2\widehat\sigma_{g,M}^2}{\widehat A_M^2-1}.
\tag{E.41}
\] The sample mean at an exact root is zero, so the variance expression
in (E.41) is the centered empirical variance. Let \(\widehat H_M\) be
the normalized conjugacy of the estimated field. Root-selection failures
and violations of the required geometry must be reported as inapplicable
predictions, rather than removed from a success denominator.

For unconditional weak-convergence and expectation statements, complete
otherwise undefined forecasts on the failure event by an arbitrary fixed
valid joint exit law, for example exit at time one with equiprobable
directions. This mathematical convention makes the random forecasts
defined everywhere. Since the failure probability tends to zero, it does
not change the limits below; it does not reclassify failed calibrations
as successful predictions. Finite-sample bounds remain restricted to
their stated good-calibration event.

\subsubsection{Proposition E.5. Field concentration, root concentration,
and root
CLT}\label{calibration-proposition-e.5.-field-concentration-root-concentration-and-root-clt}

For \(M\ge2\), \(0<\alpha<1/2\), write \[
d_M(\alpha)=\sqrt{\frac{\log(C_0M/\alpha)}{M}},
\qquad
r_M(\alpha)=\sqrt{\frac{\log(C_0/\alpha)}{M}},
\tag{E.42}
\] where \(C_0\) is a sufficiently large fixed constant.

With probability at least \(1-\alpha\), \[
\|\widehat v_M-v\|_{C^3(\mathcal K)}
\le C d_M(\alpha).
\tag{E.43}
\] For \(M\) sufficiently large that the right-hand side is below the
perturbation threshold of Lemma E.4, the same event can be chosen so
that \[
|\widehat q_M-q|\le C r_M(\alpha),
\qquad
|\widehat\sigma_{g,M}^2-\sigma_g^2|\le C r_M(\alpha),
\tag{E.44}
\] \[
|\widehat s_M-s|+|\log\widehat A_M-\log A|
\le C d_M(\alpha).
\tag{E.45}
\] Also, \[
\sqrt M(\widehat q_M-q)
=-\frac1{\lambda\sqrt M}\sum_{i=1}^M U(q,\widetilde Z_i)
+o_{\mathbb P}(1)
\Rightarrow N(0,\sigma_g^2/\lambda^2).
\tag{E.46}
\] Statements about \(\widehat q_M\) are unaffected by assigning
arbitrary values on the root-selection failure event, whose probability
tends to zero.

\textbf{Proof.} For derivative orders \(j=0,1,2,3\), the variables
\(\partial_x^j U(x,\widetilde Z_i)\) are uniformly bounded. Their sample
averages and expectations are uniformly Lipschitz in \(x\), using the
bounded derivative of order \(j+1\). Take a deterministic grid on
\(\mathcal K\) with mesh at most \(1/M\); it contains at most
\(C_{\mathcal K}M\) points. Hoeffding's inequality, a union bound over
these points and four derivative orders, and an \(O(1/M)\) interpolation
error prove (E.43).

On this event, Lemma E.4 selects a unique root near \(q\), and the
sample derivative on the segment joining the roots is at least a fixed
positive fraction of \(\lambda\). The exact root equation gives \[
\widehat q_M-q
=-\frac{\widehat v_M(q)}
{\int_0^1\widehat v_M'(q+u(\widehat q_M-q))\,du}.
\tag{E.47}
\] A separate Hoeffding bound at the single point \(q\) gives the first
inequality in (E.44), without the grid logarithm. A pointwise bound for
the average of \(U(q,\widetilde Z_i)^2\), together with \[
|U(x,z)^2-U(q,z)^2|\le2B_UL_U|x-q|,
\] gives the variance bound. Lemma E.4 and smoothness of the expression
in (E.41), with \(\lambda,\sigma_g^2\) bounded away from zero in the
fixed neighborhood, give (E.45). Split the confidence budget among these
events to obtain the stated probability.

For the CLT, the denominator in (E.47) converges in probability to
\(\lambda\). The numerator, multiplied by \(\sqrt M\), converges to
\(N(0,\sigma_g^2)\). This proves (E.46). It is the usual
estimating-equation expansion {[}E-R1{]}, applied here to the entropy
velocity. \(\square\)

\subsubsection{Theorem E.6. Prediction error for real SGD with
independently estimated
geometry}\label{calibration-theorem-e.6.-prediction-error-for-real-sgd-with-independently-estimated-geometry}

Assume G1--G4. All estimated quantities below are measurable with
respect to a calibration sigma-field \(\mathcal C_M\) independent of the
actual adaptation sample array. The initialization \(x_{0,n}\) remains
the original deterministic initialization.

Let \(\widehat a_n,\widehat c_n\) be calibrated risk or interval
endpoints. On a calibration event \(\mathcal E\), assume that the
estimated field satisfies Lemma E.4, that \[
d_b=\max(|\widehat a_n-a_n|,|\widehat c_n-c_n|)
\] is sufficiently small, and that \(\widehat\sigma_g^2>0\) is
sufficiently close to \(\sigma_g^2\). Then the estimated root is inside
the estimated interval with positive boundary distances. Define \[
\widehat\mu_n=\sqrt n\,\widehat H(x_{0,n}),\qquad
\widehat B_n^+=\widehat H(\widehat c_n),\qquad
\widehat B_n^-=-\widehat H(\widehat a_n),
\] \[
\widehat\Psi_n^\pm(t)
=\Phi\!\left(
\frac{\pm\widehat\mu_n-\sqrt n\,\widehat A^{-t}
\widehat B_n^\pm}{\widehat s}
\right).
\tag{E.48}
\] These formulas predict the true exit from \(I_n\); they do not
replace the true process by SGD driven by the estimated field.

On \(\mathcal E\), uniformly over integer \(t\ge0\), \[
\max_\pm\left|
\Pr(\tau_n\le t,J_n=\pm1\mid\mathcal C_M)
-\widehat\Psi_n^\pm(t)
\right|
\le Ce_n+C\mathcal D_{M,n},
\tag{E.49}
\] where \[
\mathcal D_{M,n}
=|\widehat\mu_n-\mu_n|
+|\widehat s-s|
+\max_\pm|\log\widehat B_n^\pm-\log B_n^\pm|
+(1+\log n)|\log\widehat A-\log A|.
\tag{E.50}
\] Let \(d_q=\widehat q-q\), \(d_v=\|\widehat v-v\|_{C^3(\mathcal K)}\),
and \(d_\sigma=|\widehat\sigma_g^2-\sigma_g^2|\). A bound in terms of
the estimated inputs is \[
\sup_{t\ge0}\max_\pm\left|
\Pr(\tau_n\le t,J_n=\pm1\mid\mathcal C_M)
-\widehat\Psi_n^\pm(t)
\right|
\le
C\left[
e_n+\sqrt n|d_q|+\sqrt n d_q^2
+(1+\log n)d_v+d_\sigma+d_b
\right].
\tag{E.51}
\] The right sides may be truncated at one.

For the empirical calibration estimator of Proposition E.5, this yields,
on its probability-\((1-\alpha)\) event and any simultaneous endpoint
event, \[
\sup_{t\ge0}\max_\pm\left|
\Pr(\tau_n\le t,J_n=\pm1\mid\mathcal C_M)
-\widehat\Psi_n^\pm(t)
\right|
\le C\left[
e_n+\sqrt n\,r_M(\alpha)
+\sqrt n\,r_M(\alpha)^2
+(1+\log n)d_M(\alpha)+d_b
\right].
\tag{E.52}
\] An endpoint-event failure probability must be added separately;
Proposition E.5 does not estimate risk endpoints.

\textbf{Proof.} Independence and the fixed initialization imply \[
\Pr(\tau_n\le t,J_n=\pm1\mid\mathcal C_M)
=\Pr(\tau_n\le t,J_n=\pm1).
\tag{E.53}
\] Apply Theorem E.3 to the right side. It remains to compare two
deterministic Gaussian formulas, treating a realized calibration sample
as fixed.

For \(B>0,a>0,s>0\), put \[
f_t(\mu,s,B,a)
=\Phi\!\left(\frac{\mu-\sqrt n B e^{-at}}s\right).
\] On the stated neighborhood, \(s,B,a\) are bounded above and away from
zero. The true \(\mu_n\) is bounded. Changing \(\mu\) first costs at
most \(C|\widehat\mu-\mu|\), with no bound needed on \(\widehat\mu\).
With \(\mu\) thereafter fixed at its true value, the derivatives with
respect to \(s\) and \(\log B\) are bounded uniformly in \(t,n\). The
latter uses \[
\sup_{y\ge0} y\,\varphi((\mu-y)/s)<\infty.
\] For \(a=\log A\), the derivative has magnitude \[
\frac{t y}{s}\varphi((\mu-y)/s),
\qquad y=\sqrt n B e^{-at}.
\] Since \[
t=\frac{\log(\sqrt n B)-\log y}{a},
\] and both \(y\varphi((\mu-y)/s)\) and \(y|\log y|\varphi((\mu-y)/s)\)
are uniformly bounded, this derivative is at most \(C(1+\log n)\). The
same estimates hold for the negative branch and along the parameter
interpolation. The mean-value bound proves (E.49).

Lemma E.4 gives \[
|\widehat\mu_n-\mu_n|
\le\sqrt n|d_q|+C(d_v+\sqrt n d_q^2),
\] \[
\max_\pm|\widehat B_n^\pm-B_n^\pm|\le C(d_v+d_b),
\qquad |\log\widehat A-\log A|\le Cd_v.
\] The boundary coordinates are bounded away from zero, so their
logarithms satisfy the same local bound. Finally, \[
|\widehat s-s|\le C(d_v+d_\sigma).
\] This proves (E.51), and substitution of Proposition E.5 proves
(E.52). \(\square\)

\textbf{Interpretation.} Peak estimation shifts a seed multiplied by
\(\sqrt n\). Growth-rate error accumulates over a logarithmic number of
updates. Field-shape and risk-endpoint errors change the transformed
distance to the boundary. Equation (E.49) is a bound for the prediction
of the actual nonlinear SGD, not a bound comparing two artificial
Gaussian processes.

\textbf{Consistency consequence.} With exact endpoints, or with
\(d_b\to_{\mathbb P}0\), empirical field calibration from the true \(Q\)
gives \[
\sup_{t\ge0}\max_\pm
\left|\Pr(\tau_n\le t,J_n=\pm1)-\widehat\Psi_n^\pm(t)\right|
\to_{\mathbb P}0
\quad\text{if }M/n\to\infty.
\tag{E.54}
\] Indeed, (E.46) gives \(d_q=O_{\mathbb P}(M^{-1/2})\), while
\(d_v=O_{\mathbb P}(\sqrt{\log M/M})\). Under \(M/n\to\infty\), every
term in (E.51) vanishes. The grid logarithm in (E.52) is not a claimed
necessary sample-complexity factor.

This conclusion is for the stated empirical-root plug-in predictor under
fixed positive \(\eta,\lambda\). It is not a minimax lower bound for all
possible predictors.

\subsection{5. A calibration limit law: consistent peaks can give
inconsistent
predictions}\label{calibration-a-calibration-limit-law-consistent-peaks-can-give-inconsistent-predictions}

The threshold scale of the frozen SGD seed is \(s/\sqrt n\). The root
CLT gives the corresponding calibration scale
\(\sigma_g/(\lambda\sqrt M)\). Their squared ratio is \[
\mathcal R_{M,n}^2
=\frac{n\sigma_g^2}{M\lambda^2s^2}
=\frac nM\frac{2+\eta\lambda}{\eta\lambda}.
\tag{E.55}
\] The cancellation of \(\sigma_g^2\) follows from comparing two
uncertainties generated by the same score law and the same update \(U\).
It need not hold for calibration from a different distribution.

\subsubsection{Theorem E.7. Random limits of independently calibrated
direction and time
predictions}\label{calibration-theorem-e.7.-random-limits-of-independently-calibrated-direction-and-time-predictions}

Assume G1--G4 and empirical calibration from the true \(Q\), with
\(M=M_n\to\infty\). Suppose \[
\frac{\sqrt n(x_{0,n}-q)}s\to z\in\mathbb R,\qquad
\mathcal R_{M,n}\to r\in[0,\infty).
\tag{E.56}
\] Let \(\pi_n=\Pr(J_n=+1)\), and define the calibrated direction
predictor \[
\widehat\pi_{M,n}=\Phi(\widehat\mu_n/\widehat s_M).
\tag{E.57}
\] Then \[
\pi_n\to\Phi(z),\qquad
\widehat\pi_{M,n}\Rightarrow\Phi(z-rZ),
\quad Z\sim N(0,1).
\tag{E.58}
\] The weak convergence in the second statement is over calibration
samples, not over additional noise inserted into the adaptation
recursion. For \(r>0\), this predictor is not consistent for \(\pi_n\).

For a time statement, suppose also \[
a_n\to a,\qquad c_n\to c,\qquad a<q<c,\qquad
\max(|\widehat a_n-a_n|,|\widehat c_n-c_n|)
\to_{\mathbb P}0.
\tag{E.59}
\] Write \(B^+=H(c)\), \(B^-=-H(a)\). Along any subsequence of batch
sizes, still indexed by \(n\), and integer times \(t_n\) satisfying \[
\sqrt n A^{-t_n}\to\zeta\in(0,\infty),
\tag{E.60}
\] the true joint CDFs satisfy \[
\Pr(\tau_n\le t_n,J_n=+1)\to\Phi(z-\zeta B^+/s),
\] \[
\Pr(\tau_n\le t_n,J_n=-1)\to\Phi(-z-\zeta B^-/s).
\tag{E.61}
\] The corresponding calibrated predictions converge jointly in
distribution to \[
\left(
\Phi(z-rZ-\zeta B^+/s),
\ \Phi(-z+rZ-\zeta B^-/s)
\right).
\tag{E.62}
\] The same calibration normal \(Z\) appears with opposite signs in the
two branches. The integer sequence condition (E.60) retains the
discrete-time phase.

\textbf{Proof.} Equation (E.46) and (E.55) give \[
\frac{\sqrt n(\widehat q_M-q)}s\Rightarrow rZ.
\tag{E.63}
\] Finite \(r\) implies \(M\) is at least of order \(n\) eventually.
Thus \(d_v\to_{\mathbb P}0\), \(\sqrt n d_q^2\to_{\mathbb P}0\), and
Lemma E.4 implies \[
\frac{\widehat\mu_n-\mu_n}s
=-\frac{\sqrt n(\widehat q_M-q)}s+o_{\mathbb P}(1).
\tag{E.64}
\] Also \(\mu_n/s\to z\) and \(\widehat s_M/s\to_{\mathbb P}1\). Theorem
E.3 and Slutsky's theorem prove (E.58).

For the time statement, \(t_n=O(\log n)\), and Proposition E.5 implies
\[
t_n|\log\widehat A_M-\log A|\to_{\mathbb P}0.
\] Hence \(\widehat A_M^{-t_n}/A^{-t_n}\to_{\mathbb P}1\). Lemma E.4 and
(E.59) give \(\widehat B_n^\pm\to_{\mathbb P}B^\pm\). Substitute these
facts, (E.63), and (E.64) into the true and calibrated formulas. Joint
convergence follows because both calibrated branches use the same
estimated root. \(\square\)

\textbf{Local bias variant.} The same proof applies if a different
calibration scheme satisfies \[
\frac{\sqrt n(\widehat q-q)}s\Rightarrow b+rZ,
\tag{E.65}
\] while its other geometry estimates have the convergence and
growth-rate control used above. Replace \(z-rZ\) by \(z-b-rZ\) in the
positive branch and \(-z+rZ\) by \(-z+b+rZ\) in the negative branch. A
fixed nonzero source-to-target peak bias does not satisfy this
finite-\(b\) assumption; \href{risk_boundaries.md}{Section 6} treats
that distinct regime.

\subsubsection{Corollary E.8. Explicit direction-prediction error at
critical
centering}\label{calibration-corollary-e.8.-explicit-direction-prediction-error-at-critical-centering}

Under Theorem E.7, \[
\mathbb E_{\rm cal}\widehat\pi_{M,n}
\longrightarrow
\Phi\!\left(\frac z{\sqrt{1+r^2}}\right).
\tag{E.66}
\] At \(z=0\), \[
\mathbb E_{\rm cal}|\widehat\pi_{M,n}-\pi_n|
\longrightarrow\frac{\arctan r}{\pi}.
\tag{E.67}
\] In particular, \(r=1\) gives limiting mean absolute probability error
\(1/4\), even though \(\widehat q_M\to q\). At this value the limiting
predictor is uniform on \((0,1)\).

\textbf{Proof.} All predictors lie in \([0,1]\), so weak convergence in
(E.58) also gives convergence of their bounded continuous moments. For
an independent standard normal \(Z_0\), \[
\mathbb E\Phi(z-rZ)
=\Pr(Z_0+rZ\le z)
=\Phi(z/\sqrt{1+r^2}).
\] At \(z=0\), symmetry gives \[
f(r):=\mathbb E|\Phi(rZ)-1/2|
=2\int_0^\infty\{\Phi(ru)-1/2\}\varphi(u)\,du.
\] Differentiation under the integral yields \[
f'(r)=2\int_0^\infty u\varphi(ru)\varphi(u)\,du
=\frac1{\pi(1+r^2)}.
\] Since \(f(0)=0\), (E.67) follows. For \(r=1\), \(\Phi(-Z)\) is
uniform. \(\square\)

Equation (E.66) has the algebraic appearance of an inflated seed
variance, but it is an average of random forecasts. The actual direction
probability remains \(\Phi(z)\). Replacing the actual seed variance by
\(s^2(1+r^2)\) would therefore answer a different question.

If \(M\to\infty\) but \(M/n\to0\), the same empirical-root predictor
instead converges in distribution to a variable taking the values zero
and one with probabilities \(1/2,1/2\), for every fixed finite \(z\). To
verify this, apply (E.46) on the calibration scale:
\(\sqrt M(x_{0,n}-q)\to0\), so the sign of \(x_{0,n}-\widehat q_M\) is
asymptotically balanced. Its magnitude multiplied by \(\sqrt n\)
diverges in probability away from an arbitrarily small neighborhood of
zero. Monotonicity and the local derivative normalization of
\(\widehat H_M\) preserve this sign and growth. The predictor becomes
extreme across calibration repetitions; it does not become accurate.

{\def\LTcaptype{table} % do not increment counter
\begin{longtable}[]{@{}
  >{\raggedright\arraybackslash}p{(\linewidth - 2\tabcolsep) * \real{0.5000}}
  >{\raggedright\arraybackslash}p{(\linewidth - 2\tabcolsep) * \real{0.5000}}@{}}
\toprule\noalign{}
\begin{minipage}[b]{\linewidth}\raggedright
Relative calibration scale
\end{minipage} & \begin{minipage}[b]{\linewidth}\raggedright
Direction prediction in the critical regime
\end{minipage} \\
\midrule\noalign{}
\endhead
\bottomrule\noalign{}
\endlastfoot
\(\mathcal R_{M,n}\to0\) & Consistent for the actual direction
probability \\
\(\mathcal R_{M,n}\to r\in(0,\infty)\) & Nonvanishing random prediction
error \\
\(\mathcal R_{M,n}\to\infty\), with \(M\to\infty\) & Random extreme
predictions for the empirical-root estimator \\
\end{longtable}
}

These statements fix \(Q,w,\eta,\lambda\). The factor in (E.55) can be
large when \(\eta\lambda\) is small, but the results do not establish a
joint limit with \(\eta\to0\) or \(\lambda\to0\).

\subsection{References}\label{calibration-references}

{[}E-R1{]} A. W. van der Vaart. \emph{Asymptotic Statistics}. Cambridge
University Press, 1998. Chapter 5, ``M- and Z-Estimators.''
\href{https://doi.org/10.1017/CBO9780511802256.006}{Publisher chapter}.

{[}E-R2{]} Y. Bakhtin. \emph{Exit asymptotics for small diffusion about
an unstable equilibrium}. arXiv:math/0701569, 2007.
\href{https://arxiv.org/abs/math/0701569}{Preprint}. Journal DOI:
\href{https://doi.org/10.1016/j.spa.2007.06.003}{10.1016/j.spa.2007.06.003}.

{[}E-R3{]} Y. Bakhtin and A. Gaál. \emph{Exit time asymptotics for
dynamical systems with fast random switching near an unstable
equilibrium}. arXiv:1901.05513, 2019.
\href{https://arxiv.org/abs/1901.05513}{Preprint}. Journal DOI:
\href{https://doi.org/10.1142/S0219493720500264}{10.1142/S0219493720500264}.

{[}E-R4{]} P. Massart. \emph{The Tight Constant in the
Dvoretzky--Kiefer--Wolfowitz Inequality}. The Annals of Probability
18(3):1269--1283, 1990.
\href{https://doi.org/10.1214/aop/1176990746}{DOI}.
